% RQ1 and RQ3: the multi-agent learning study. Numbers from numbers.tex / tab_main.tex.
\subsection{RQ1: violations during learning}
\label{sec:rq1}
\Cref{fig:training} tracks certified coverage and cumulative violations over training.
\begin{itemize}[leftmargin=1.2em,itemsep=1pt,topsep=1pt]
\item \emph{Projection in the environment.} It executed \emph{zero} infeasible assignments
  over $\nOursTrainSteps$ agent-epochs, pooled across \nNSeeds\ seeds, from the first update
  onwards. Per-agent masking also reached zero, as separability predicts. For both methods
  zero holds by construction, so the CP bound in \cref{tab:main} only quantifies the
  exposure observed. That bound assumes independent agent-epochs, which these are not.
\item \emph{PPO-Lagrangian.} Although it observes the certificate, it executed
  \nLagTrainViol\ infeasible assignments while learning, $\nLagTrainRate\%$ of agent-epochs.
  Its multiplier then rose until the policy largely stopped serving tasks. Final certified
  coverage is $\nLagCov$, and evaluation violations still average $\nLagViolMean\%$ of
  agent-epochs, above the 0.5\% limit.
\item \emph{Jamming-blind PPO.} It executed \nBlindTrainViol\ infeasible assignments
  ($\nBlindTrainRate\%$).
\item \emph{Stronger in-expectation learners} (\cref{tab:r2}). MAPPO-Lagrangian with a cost
  critic executed \nMLagTrainViol\ infeasible assignments while learning
  ($\nMLagTrainRate\%$), and \nMLagZeroTrainViol\ with a cost limit of zero.
  PID-Lagrangian executed \nPidLagTrainViol\ and then stopped serving tasks altogether.
  QMIX without masks executed \nQmixTrainViol. Only the two methods that remove infeasible
  actions by construction, our layer and masking (with PPO or with QMIX), learned without a
  violation.
\end{itemize}
In-expectation methods violate the constraint while they learn, and learning is when a
fleet would be flight-tested. No setting of the Lagrangian's cost limit or dual step changed
this (\cref{app:sens}).

\begin{figure*}[t]
\centering
\begin{tikzpicture}
\begin{groupplot}[group style={group size=2 by 1, horizontal sep=1.4cm},
  width=0.46\textwidth, height=4.2cm, xlabel={fleet-epochs ($\times10^3$)},
  grid=major, grid style={black!10}, tick label style={font=\scriptsize},
  label style={font=\small}, xmin=0, xmax=480,
  legend style={font=\scriptsize, fill=white, fill opacity=0.85, text opacity=1}]
\nextgroupplot[ylabel={certified coverage $r_{\mathrm{cert}}$}, ymin=0, ymax=0.8,
  legend style={at={(0.98,0.98)}, anchor=north east, font=\tiny}, legend columns=2]
\addplot[name path=ou, draw=none, forget plot] table[x=steps,y=hi]{figdata/curve_ours.dat};
\addplot[name path=ol, draw=none, forget plot] table[x=steps,y=lo]{figdata/curve_ours.dat};
\addplot[red!20, forget plot] fill between[of=ou and ol];
\addplot[red!70!black, thick] table[x=steps,y=rcert]{figdata/curve_ours.dat};
\addplot[orange!80!black, thick] table[x=steps,y=rcert]{figdata/curve_mask.dat};
\addplot[blue!70!black, thick] table[x=steps,y=rcert]{figdata/curve_lagrangian.dat};
\addplot[gray, thick, dashed] table[x=steps,y=rcert]{figdata/curve_blind.dat};
\legend{projection (ours), masking, PPO-Lagrangian, jamming-blind}
\nextgroupplot[ylabel={cumulative violations (mean per run)}, ymode=log, ymin=0.8, ymax=3e6,
  legend pos=south east]
\addplot[red!70!black, thick] table[x=steps,y=cumviol]{figdata/curve_ours.dat};
\addplot[orange!80!black, thick, densely dotted] table[x=steps,y=cumviol]{figdata/curve_mask.dat};
\addplot[blue!70!black, thick] table[x=steps,y=cumviol]{figdata/curve_lagrangian.dat};
\addplot[gray, thick, dashed] table[x=steps,y=cumviol]{figdata/curve_blind.dat};
\legend{projection (0), masking (0), PPO-Lagrangian, jamming-blind}
\end{groupplot}
\end{tikzpicture}
\caption{RQ1. Left: certified coverage during training (stochastic policy; mean over
\nNSeeds\ seeds, 10-update moving average; band: min--max across seeds for ours). Right: cumulative executed
violations on a log axis. Zero is drawn at 1 for visibility. The projection and masking
curves lie on the floor throughout.}
\label{fig:training}
\end{figure*}

\subsection{RQ3: what the policy learns and what the mechanism enforces}
\label{sec:rq3}
\Cref{tab:main} evaluates every method greedily on held-out missions. Four findings stand out.

\emph{The mechanism does most of the work.} Adding the projection at deployment to \emph{any}
learned policy removed every violation and lifted the coverage ratio to at least
$\nLagProjRatio$:
\begin{itemize}[leftmargin=1.2em,itemsep=0pt,topsep=1pt]
\item jamming-blind: certified coverage $\nBlindCov\to\nBlindProjCov$;
\item masking: $\nMaskCov\to\nMaskProjCov$;
\item PPO-Lagrangian: $\nLagCov\to\nLagProjCov$.
\end{itemize}
Without it, the jamming-blind policy and the no-certificate heuristic reach seemingly
competitive certified coverage ($\nBlindCov$ and $\nHeurNoneCov$). They do so only by holding
tasks without certification on $\nBlindViolMean\%$ and $\nHeurNoneViolMean\%$ of
agent-epochs. The jamming-blind policy in fact converges to the keep-or-nearest rule, with no
reassignments.

\emph{Masking is safe but not maximal.} Its coverage ratio of $\nMaskRatio$ is the gap that
\cref{thm:proj}(ii) closes. Agents that each pick an eligible task still collide.

\emph{Learning adds little on top of the guarantee, and where the layer sits during training
does not matter.}
\begin{itemize}[leftmargin=1.2em,itemsep=0pt,topsep=1pt]
\item The policy trained with the projection reaches $\nOursCov$ [$\nOursCovLo$,
  $\nOursCovHi$] at a coverage ratio of $\nOursRatio$, over \nNSeeds\ training seeds.
\item The keep-or-nearest heuristic with the same projection reaches $\nHeurCov$; the
  intervals only just separate.
\item Policies trained with masking, or blind, and projected at deployment do as well
  ($\nMaskProjCov$, $\nBlindProjCov$), and the deviation penalty makes no measurable
  difference ($\nNoDevCov$ without it).
\end{itemize}
What training \emph{with} the projection buys is the zero-violation learning of RQ1, not a
better final policy.

\emph{Static and design-time allocators fail in opposite ways.} The tightest-fit greedy
re-solves from scratch every epoch, keeps its aircraft in transit ($\nGnpSwitch\%$
reassignments) and certifies only $\nGnpCov$. A fixed design margin certifies $\nFixedCov$
and still violates on $\nFixedViolMean\%$ of agent-epochs, because a margin fixed at design
time cannot track a certificate that moves.

\begin{table*}[t]
\centering\footnotesize
\setlength{\tabcolsep}{3.5pt}
\caption{RQ3. Held-out evaluation ($N=12$, $M=8$; \nNSeeds\ training seeds $\times$ 3
evaluation seeds $\times$ 16 fleets for learned methods; 3 $\times$ 16 for the references
below the line). Values are IQMs with 95\% stratified-bootstrap CIs. Cov.\ ratio: feasible
holders divided by $(M-d)$. Violations are in \% of agent-epochs (IQM over runs). ``$0\,(\le
x)$'': none observed, with the 95\% CP upper bound on the rate, assuming independent
agent-epochs. Train viol.: executed infeasible assignments over all of training.}
\label{tab:main}
\resizebox{\linewidth}{!}{%
\begin{tabular}{@{}lcccccc@{}}
\toprule
Method & $r_{\mathrm{cert}}$ & cov.\ ratio & eval viol. [\%] & train viol. & exposure [\%] & reassign. [\%]\\
\midrule
\textbf{Projection (SE) + PPO (ours)} & $0.496$ {\scriptsize[0.473, 0.525]} & $1.00$ & $0$ ($\le3.5\times10^{-6}$) & $0$ ($\le1.0\times10^{-7}$) & $1.62$ & $4.9$ \\
\quad ablation: no deviation penalty & $0.499$ {\scriptsize[0.484, 0.520]} & $1.00$ & $0$ ($\le3.5\times10^{-6}$) & $0$ ($\le1.0\times10^{-7}$) & $1.66$ & $4.9$ \\
Per-agent action masking + PPO & $0.464$ {\scriptsize[0.458, 0.473]} & $0.62$ & $0$ ($\le3.5\times10^{-6}$) & $0$ ($\le1.0\times10^{-7}$) & $0.52$ & $0.7$ \\
\quad + projection at deployment & $0.499$ {\scriptsize[0.489, 0.508]} & $1.00$ & $0$ ($\le3.5\times10^{-6}$) & -- & $1.71$ & $4.9$ \\
PPO-Lagrangian (certificate observed) & $0.036$ {\scriptsize[0.028, 0.047]} & $0.04$ & $0.02$ & $3{,}903{,}179$ & $0.03$ & $0.0$ \\
\quad + projection at deployment & $0.393$ {\scriptsize[0.383, 0.405]} & $0.99$ & $0$ ($\le3.5\times10^{-6}$) & -- & $1.48$ & $3.9$ \\
Jamming-blind PPO & $0.451$ {\scriptsize[0.439, 0.467]} & $0.51$ & $13.75$ & $4{,}744{,}505$ & $0.37$ & $0.0$ \\
\quad + projection at deployment & $0.492$ {\scriptsize[0.478, 0.508]} & $1.00$ & $0$ ($\le3.5\times10^{-6}$) & -- & $1.71$ & $5.0$ \\
\midrule
Keep-or-nearest + projection & $0.435$ {\scriptsize[0.395, 0.477]} & $1.00$ & $0$ ($\le1.7\times10^{-5}$) & -- & $1.45$ & $2.3$ \\
Tightest-fit greedy, re-solved each epoch & $0.169$ {\scriptsize[0.142, 0.216]} & $0.98$ & $0$ ($\le1.7\times10^{-5}$) & -- & $1.99$ & $8.6$ \\
Fixed design margin (1\,s dead reckoning) & $0.163$ {\scriptsize[0.155, 0.170]} & $0.31$ & $8.03$ & -- & $0.20$ & $0.0$ \\
Keep-or-nearest, no certificate & $0.459$ {\scriptsize[0.433, 0.484]} & $0.51$ & $13.97$ & -- & $0.38$ & $0.0$ \\
\bottomrule
\end{tabular}}

\end{table*}

\subsection{RQ3 (continued): stronger baselines and a classical auction}
\label{sec:rq3b}
\Cref{tab:r2} adds the stronger learners and a non-learned consensus auction at the
headline size. Three findings matter for the security argument.

\emph{Better in-expectation learners still violate, and the monitor still repairs them.}
MAPPO-Lagrangian with a cost critic is the strongest unconstrained learner: it reaches
$\nMLagCov$ certified coverage and, once trained, violates on only $\nMLagEvalViolPct\%$ of
agent-epochs. Its guarantee is nonetheless in expectation: it paid for that policy with
\nMLagTrainViol\ violations during learning, and lowering its cost limit to zero did not
remove them. With the projection at deployment, every learned policy in the table executes
no violation and reaches a coverage ratio of at least $\nProjMinRatio$.

\emph{A classical auction under the same constraint beats every learned allocator.} CBBA
with feasibility enforced inside the bundle certifies $\nCbbaCov$ [$\nCbbaCovLo$,
$\nCbbaCovHi$] of the task value, against $\nOursCov$ [$\nOursCovLo$, $\nOursCovHi$] for the
best learned policy, with no violation and a coverage ratio of $\nCbbaRatio$. The comparison
favours the auction in one respect: its consensus phase gives every aircraft the fleet's full
state at every epoch, while the learned policies act on local observations in one shot. In
our setting, that information is worth more than learning.

\emph{The result supports the paper's claim, not a learning claim.} Our guarantees do not
depend on which allocator proposes: they hold for the auction, for the heuristic, for every
learned policy and for an adversarially chosen one (A4). The learning study shows that the
monitor can sit in the training loop at no cost to learning; it does not show that learning
is the best allocator, and we do not claim it.

\begin{table}[t]
\centering\footnotesize
\setlength{\tabcolsep}{2.5pt}
\caption{RQ3. Stronger baselines at $N=12$ (\nMLagSeeds\ training seeds $\times$ 3
evaluation seeds $\times$ 16 fleets; CBBA is deterministic, 3 $\times$ 16). Columns as in
\cref{tab:main}; eval.\ violations in \% of agent-epochs.}
\label{tab:r2}
\resizebox{\columnwidth}{!}{\begin{tabular}{@{}lcccc@{}}
\toprule
Method & $r_{\mathrm{cert}}$ & cov.\ ratio & eval viol. [\%] & train viol.\\
\midrule
\textbf{Projection (SE) + PPO (ours)} & $0.491$ {\scriptsize[0.477, 0.509]} & $1.00$ & $0$ ($\le1.7\times10^{-6}$) & $0$ ($\le5.2\times10^{-8}$) \\
Keep-or-nearest + projection & $0.435$ {\scriptsize[0.395, 0.477]} & $1.00$ & $0$ ($\le1.7\times10^{-5}$) & -- \\
\midrule
MAPPO-Lagrangian, cost critic, limit $0.005$ & $0.445$ {\scriptsize[0.438, 0.451]} & $0.57$ & $0.027$ & $2{,}897{,}602$ \\
\quad + projection at deployment & $0.463$ {\scriptsize[0.455, 0.473]} & $0.99$ & $0$ ($\le1.7\times10^{-6}$) & -- \\
MAPPO-Lagrangian, cost critic, limit $0$ & $0.442$ {\scriptsize[0.436, 0.451]} & $0.57$ & $0.091$ & $2{,}929{,}578$ \\
\quad + projection at deployment & $0.466$ {\scriptsize[0.455, 0.477]} & $1.00$ & $0$ ($\le1.7\times10^{-6}$) & -- \\
PID-Lagrangian PPO & $0.000$ {\scriptsize[0.000, 0.000]} & $0.00$ & $0.040$ & $8{,}089{,}556$ \\
\quad + projection at deployment & $0.347$ {\scriptsize[0.336, 0.360]} & $0.98$ & $0$ ($\le1.7\times10^{-6}$) & -- \\
QMIX (certificate observed) & $0.422$ {\scriptsize[0.410, 0.436]} & $0.60$ & $8.473$ & $7{,}526{,}523$ \\
\quad + projection at deployment & $0.460$ {\scriptsize[0.446, 0.477]} & $1.00$ & $0$ ($\le1.7\times10^{-6}$) & -- \\
QMIX with eligibility masks & $0.451$ {\scriptsize[0.444, 0.459]} & $0.68$ & $0$ ($\le1.7\times10^{-6}$) & $0$ ($\le5.2\times10^{-8}$) \\
\quad + projection at deployment & $0.455$ {\scriptsize[0.437, 0.473]} & $1.00$ & $0$ ($\le1.7\times10^{-6}$) & -- \\
CBBA, feasibility in the bundle & $0.539$ {\scriptsize[0.525, 0.561]} & $1.00$ & $0$ ($\le1.7\times10^{-5}$) & -- \\
\bottomrule
\end{tabular}
}
\end{table}

\subsection{RQ3 (continued): fleet size}
\label{sec:rq3c}
\Cref{tab:scale,fig:scale} repeat the core methods at $N\in\{\nScaleSizes\}$.
\begin{itemize}[leftmargin=1.2em,itemsep=1pt,topsep=1pt]
\item \emph{The guarantees scale.} Our layer and masking executed no infeasible assignment
  during training at any size, and our coverage ratio was at least $\nScaleMinRatio$
  throughout. PPO-Lagrangian executed millions of violations at every size and, from
  $N=24$ on, converged to idling ($\nScaleTwentyFourLag$ certified coverage).
\item \emph{Learning does not scale with our budget.} Certified coverage of the learned
  policy fell from $\nScaleTwelveOurs$ at $N=12$ to $\nScaleFortyEightOurs$ at $N=48$, with
  masking plus projection matching it at every size. It stayed above keep-or-nearest plus
  projection, separated at $N=\nScaleBeatsHeur$ and within the confidence intervals elsewhere.
\item \emph{The auction does.} CBBA certified $\nScaleTwelveCbba$, $\nScaleTwentyFourCbba$ and
  $\nScaleFortyEightCbba$ of the task value at the three sizes.
\end{itemize}
Our budget per size was 200 updates on 4 CPU cores, below the runbook's 400--600; larger
budgets may narrow the gap to the auction, and the artifact runs them unchanged.

\begin{table*}[t]
\centering\footnotesize
\caption{RQ3, fleet size. Certified coverage (IQM, 95\% stratified-bootstrap CI) of the
learned policy with the projection in the environment, masking with the projection at
deployment, and two non-learned allocators; $M=\lceil2N/3\rceil$, arena and jammers scaled
with $N$.}
\label{tab:scale}
\begin{tabular}{@{}rrcccccc@{}}
\toprule
$N$ & seeds & ours $r_{\mathrm{cert}}$ & cov.\ ratio & masking + proj. & keep-or-nearest + proj. & CBBA & Lagrangian train viol.\\
\midrule
12 & 10 & $0.491$ {\scriptsize[0.477, 0.509]} & $1.00$ & $0.498$ {\scriptsize[0.492, 0.505]} & $0.435$ {\scriptsize[0.395, 0.477]} & $0.539$ {\scriptsize[0.525, 0.561]} & $8{,}193{,}222$ \\
24 & 5 & $0.481$ {\scriptsize[0.446, 0.522]} & $1.00$ & $0.474$ {\scriptsize[0.437, 0.516]} & $0.377$ {\scriptsize[0.343, 0.439]} & $0.526$ {\scriptsize[0.446, 0.588]} & $4{,}583{,}986$ \\
48 & 5 & $0.386$ {\scriptsize[0.350, 0.433]} & $1.00$ & $0.387$ {\scriptsize[0.353, 0.439]} & $0.316$ {\scriptsize[0.246, 0.401]} & $0.535$ {\scriptsize[0.474, 0.580]} & $4{,}685{,}248$ \\
96 & 3 & $0.009$ {\scriptsize[0.008, 0.011]} & $1.00$ & $0.082$ {\scriptsize[0.077, 0.088]} & $0.272$ {\scriptsize[0.197, 0.316]} & $0.473$ {\scriptsize[0.416, 0.543]} & $1{,}976{,}207$ \\
\bottomrule
\end{tabular}

\end{table*}

\begin{figure}[t]
\centering
\begin{tikzpicture}
\begin{axis}[width=\columnwidth, height=4.6cm, xmode=log, log basis x=2,
  xtick={12,24,48,96}, xticklabels={12,24,48,96}, xlabel={fleet size $N$},
  ylabel={certified coverage $r_{\mathrm{cert}}$}, ymin=0, ymax=0.7, grid=major,
  grid style={black!10}, tick label style={font=\scriptsize}, label style={font=\small},
  legend style={font=\scriptsize, at={(0.5,-0.32)}, anchor=north, legend columns=2, draw=none}]
\addplot[red!70!black, thick, mark=*, mark size=1.4pt, error bars/.cd, y dir=both, y explicit]
  table[x=N, y=rcert, y error plus expr=\thisrow{hi}-\thisrow{rcert}, y error minus expr=\thisrow{rcert}-\thisrow{lo}]{figdata/scale_ours.dat};
\addplot[orange!80!black, thick, densely dashed, mark=square*, mark size=1.3pt] table[x=N, y=rcert]{figdata/scale_mask_proj.dat};
\addplot[gray, thick, mark=triangle*, mark size=1.5pt] table[x=N, y=rcert]{figdata/scale_heuristic_proj.dat};
\addplot[green!50!black, thick, mark=diamond*, mark size=1.6pt, error bars/.cd, y dir=both, y explicit]
  table[x=N, y=rcert, y error plus expr=\thisrow{hi}-\thisrow{rcert}, y error minus expr=\thisrow{rcert}-\thisrow{lo}]{figdata/scale_cbba.dat};
\addplot[blue!70!black, thick, dotted, mark=o, mark size=1.3pt] table[x=N, y=rcert]{figdata/scale_lagrangian.dat};
\legend{projection + PPO (ours), masking + projection, keep-or-nearest + projection, CBBA (feasible bundles), PPO-Lagrangian}
\end{axis}
\end{tikzpicture}
\caption{RQ3, fleet size. Every method shown except PPO-Lagrangian executes no infeasible
assignment; the learned allocators lose coverage as the fleet grows at a fixed training
budget, while the auction does not.}
\label{fig:scale}
\end{figure}
